\chapter{Binding / Knowledge Soundness of the IPA Scheme}

\providecommand{\oSpec}{\mathcal{O}_{\mathrm{spec}}}

This chapter states the second headline security property of kimchi's IPA
polynomial commitment: \emph{knowledge soundness}, in the form of
\emph{evaluation binding}. Where the correctness chapter says that honest
openings are accepted, binding says the dual: it is computationally infeasible to
make the verifier accept two \emph{different} evaluations of the same commitment
at the same point. In the Halo development this is packaged as \emph{computational
witness-extended emulation} (Halo Definition~2, Theorem~1): for any efficient
malicious prover there is an emulator that, alongside a transcript indistinguishable
from the real interaction, extracts a witness satisfying the opening relation
(\cref{def:ipa_relation}) whenever the transcript accepts. We record Halo's
definition verbatim, then state the property in the ArkLib experiment shape
(\cref{def:commitment_binding}) that the Lean development targets, and give an
honest high-level sketch of the extractor argument. The full extraction proof is
Halo Appendix~A and is deferred to a later research phase; the sketch here is a
sketch, not a proof.

\section{Witness-extended emulation}

\begin{definition}[Computational witness-extended emulation]
  \label{def:wee}
  \lean{Kimchi.Commitment.IPA.Binding.witnessExtendedEmulation}
  \uses{def:ipa_relation}
  % SOURCE: [Halo], Definition 2 (Computational Witness-Extended Emulation) (read from references/halo.txt)
  % SOURCE QUOTE: "Definition 2 (Computational Witness-Extended Emulation).
  % (Setup, P, V) has witness-extended emulation if for all deterministic polynomial-
  % time P ∗ there exists an expected polynomial-time emulator E such that for all
  % pairs of interactive adversaries A1 , A2
  %
  %                                         σ ← Setup(1λ )
  %
  %
  %       Pr                   A1 (tr) = 1 (x, s) ← A2 (σ)            −
  %
  %
  %                                        tr ← hP ∗ (σ, x, s), V(σ, x)i
  %                                                                          ≤ negl(λ)
  %                                        σ ← Setup(1λ )
  %
  %                            A1 (tr) = 1
  %       Pr                              (x, s) ← A2 (σ)
  %
  %           ∧ (tr accepts ⇒ (x, w) ∈ R)
  %
  %                                        (tr, w) ← E O (σ, x)
  %
  % where the oracle is given by O = hP ∗ (σ, x, s), V(σ, x)i and permits rewinding to a
  % specific point and resuming with fresh randomness for the verifier from that point
  % onward. We also define computational witness-extended emulation by restricting
  % to non-uniform polynomial-time adversaries A1 and A2 ."
  \textit{Source: Halo, Definition~2 (computational witness-extended emulation).}
  The opening argument \((\Setup, \Open, \VerifyOpen)\) has \emph{computational
  witness-extended emulation} if for every deterministic polynomial-time malicious
  prover \(P^{*}\) there is an expected-polynomial-time emulator \(E\) such that for
  every pair of interactive adversaries \((A_{1}, A_{2})\) the two experiments
  \begin{itemize}
    \item sample \(\sigma \leftarrow \Setup(1^{\lambda})\) and \((x, s) \leftarrow
      A_{2}(\sigma)\), run the real interaction \(\mathrm{tr} \leftarrow
      \langle P^{*}(\sigma, x, s), \VerifyOpen(\sigma, x)\rangle\), and output
      \(A_{1}(\mathrm{tr})\); versus
    \item the same sampling, but with \((\mathrm{tr}, w) \leftarrow E^{O}(\sigma, x)\)
      produced by the emulator, whose oracle
      \(O = \langle P^{*}(\sigma, x, s), \VerifyOpen(\sigma, x)\rangle\) permits
      rewinding to a chosen point and resuming with fresh verifier randomness, and
      output \(A_{1}(\mathrm{tr})\) conjoined with the soundness clause
      \(\mathrm{tr}\ \text{accepts} \Rightarrow (x, w) \in R\),
  \end{itemize}
  differ in acceptance probability by at most \(\mathrm{negl}(\lambda)\). Here
  \(R\) is the opening relation of \cref{def:ipa_relation}. Thus the emulator
  reproduces the prover's transcript distribution while additionally extracting a
  witness \(w = (a, r)\) that genuinely opens the commitment whenever the
  transcript is accepting. The non-uniform variant restricts \(A_{1}, A_{2}\) to
  non-uniform polynomial time.
\end{definition}

Witness-extended emulation (\cref{def:wee}) is the \emph{intended}
knowledge-soundness notion for the IPA opening argument: it is the property whose
extractor recovers a genuine opening witness whenever a transcript accepts. The
headline result of this chapter, \cref{thm:ipa_binding}, instead establishes the
strictly weaker but fully well-formed notion of \emph{evaluation binding}
(\cref{def:commitment_binding}) --- that no efficient adversary opens one
commitment to two distinct evaluations at the same point. A faithful Lean
statement of the full witness-extended emulation property is deferred: ArkLib's
corresponding \(\extractability\) predicate currently has its body stubbed to
\(\bot\) (it has not yet been given its real definition upstream), so there is no
provable target to formalize against. The knowledge-soundness formulation is
therefore recorded here as the semantic goal but is not a Phase-1 target; it
becomes actionable only once the real ArkLib extractability body is upstreamed.

\begin{definition}[Discrete Log Relation Assumption]
  \label{def:dl_relation_assumption}
  \lean{Kimchi.Commitment.IPA.Binding.dlRelationAssumption}
  % SOURCE: [Halo], Definition 6 (Discrete Log Relation Assumption) (read from references/halo.txt)
  % SOURCE QUOTE: "Definition 6 (Discrete Log Relation Assumption). For all adversaries A
  % and for all n ≥ 2
  %         h                                                 i
  %            $
  %            − Gn ; a ∈ Fn ← A(G, G) : ∃ai 6= 0 ∧ ha, Gi = O ≤ negl(λ)
  %       Pr G ←
  %
  % The discrete log relation assumption generalizes the discrete log assumption to
  % arbitrary numbers of random group elements. We say that ha, Gi = O is a non-
  % trivial discrete log relation between elements of G when at least one of a is
  % nonzero."
  \textit{Source: Halo, Definition~6 (discrete log relation assumption).}
  Let \(G\) be a group of order the scalar field \(\F\). The \emph{discrete log
  relation assumption} postulates that for every \(n \ge 2\) and every efficient
  adversary \(A\), if \((G_{1}, \dots, G_{n}) \leftarrow G^{n}\) are sampled
  uniformly at random, then the probability that \(A\) outputs a coefficient vector
  \(a = (a_{1}, \dots, a_{n}) \in \F^{n}\) with some \(a_{i} \neq 0\) and
  \[
    \inner{a}{(G_{1}, \dots, G_{n})} = \sum_{i=1}^{n} [a_{i}]G_{i} = O
  \]
  (the group identity) is negligible in the security parameter \(\lambda\). A tuple
  \(a\) with \(\inner{a}{(G_{1}, \dots, G_{n})} = O\) and at least one \(a_{i} \neq
  0\) is called a \emph{non-trivial discrete log relation} among the generators.
  Writing \(\varepsilon_{\mathrm{dl}}(\lambda)\) for the maximal advantage of any
  efficient adversary in producing such a relation, the assumption asserts
  \(\varepsilon_{\mathrm{dl}}(\lambda)\) is negligible. This is a postulated
  hardness assumption --- a foundational cryptographic hypothesis taken as given,
  not derived from the elliptic-curve group law --- and it is the quantity the IPA
  binding error term is measured against.
\end{definition}

\begin{theorem}[Binding of the IPA scheme]
  \label{thm:ipa_binding}
  \lean{Kimchi.Commitment.IPA.Binding.ipa_binding}
  \uses{def:ipa_scheme, def:commitment_binding, def:commitment_binding_experiment,
    def:dl_relation_assumption, def:ipa_opening_protocol}
  % SOURCE: [Halo], Theorem 1 (read from references/halo.txt)
  % SOURCE QUOTE: "Theorem 1. The protocol presented in Section 3.1 has perfect completeness,
  % computational witness-extended emulation, and perfect special honest-verifier
  % zero knowledge.
  %
  % This proof appears in Appendix A."
  \textit{Source: Halo, Theorem~1 (computational witness-extended emulation of the
  Section~3.1 protocol).}
  The IPA commitment scheme \(\mathsf{ipa}\) of \cref{def:ipa_scheme} satisfies
  ArkLib evaluation \(\binding\) (\cref{def:commitment_binding}) with negligible
  error: there is a negligible \(\varepsilon_{\mathrm{b}}(\lambda)\) such that for
  every auxiliary-state type and every binding adversary, the binding experiment
  (\cref{def:commitment_binding_experiment}) for \(\mathsf{ipa}\) is at most
  \(\varepsilon_{\mathrm{b}}(\lambda)\). Concretely, no efficient adversary can
  produce a commitment \(P\), a point \(x\), and two responses \(v_{1} \neq v_{2}\)
  together with malicious openings of the protocol of
  \cref{def:ipa_opening_protocol} that both pass \(\VerifyOpen\). This is the
  ArkLib-experiment expression of Halo's computational witness-extended emulation
  (\cref{def:wee}): the binding error \(\varepsilon_{\mathrm{b}}\) is exactly the
  emulation gap of Halo Theorem~1 specialized to the evaluation query, so an
  adversary winning the binding game would refute the witness-extended emulation
  of the opening argument.
\end{theorem}

% SOURCE: [Halo], Section 3.1, Protocol Description (read from references/halo.txt)
% SOURCE QUOTE PROOF: "and that v 0 = ha, (1, x, x2 , ..., xd−1 )i. As the prover did not know U in advance,
% this establishes that v = v 0 . (In this respect we differ from the protocol described
% in [37, Appendix A.3] which effectively has U fixed prior to the argument; a
% prover with malicious control of P would then be able to interfere with the
% argument by including terms involving U in P .)"
% SOURCE: [Halo], Definition 6 (Discrete Log Relation Assumption) (read from references/halo.txt)
% SOURCE QUOTE PROOF: "Definition 6 (Discrete Log Relation Assumption). For all adversaries A
% and for all n ≥ 2
%         h                                                 i
%            $
%            − Gn ; a ∈ Fn ← A(G, G) : ∃ai 6= 0 ∧ ha, Gi = O ≤ negl(λ)
%       Pr G ←
%
% The discrete log relation assumption generalizes the discrete log assumption to
% arbitrary numbers of random group elements. We say that ha, Gi = O is a non-
% trivial discrete log relation between elements of G when at least one of a is
% nonzero."
\begin{proof}
  \uses{def:commitment_binding_adversary, def:ipa_relation, def:wee,
    def:dl_relation_assumption}
  \textit{Source: Halo, Section~3.1 and Definition~6; the full construction is Halo
  Appendix~A and is deferred to the research phase.} This is a high-level sketch,
  not a complete proof.

  Suppose, for contradiction, a binding adversary
  (\cref{def:commitment_binding_adversary}) produces a commitment \(P\), an
  evaluation point \(x\), and two distinct claims \(v_{1} \neq v_{2}\), with
  malicious provers that drive the opening protocol
  (\cref{def:ipa_opening_protocol}) to acceptance for both \((P, x, v_{1})\) and
  \((P, x, v_{2})\). Recall that the first verifier move samples a fresh group
  element \(U\) and both parties pass to \(P' = P + [v]U\); because the prover does
  not learn \(U\) before committing to \(P\), an accepting transcript pins the
  claimed value to the value actually committed (Halo \S3.1).

  Apply the emulator \(E\) of \cref{def:wee} to each accepting branch. Rewinding
  the \(k = \log_{2} d\) folding rounds with fresh challenges \(u_{j}\) and solving
  the resulting linear system across rounds extracts, for branch \(i\), a witness
  \((a_{i}, r_{i})\) with
  \[
    P = \inner{a_{i}}{g} + [r_{i}]H
    \qquad\text{and}\qquad
    v_{i} = \inner{a_{i}}{(1, x, x^{2}, \dots, x^{d-1})},
  \]
  i.e. \((P, x, v_{i}; a_{i}, r_{i})\) lies in the opening relation
  \cref{def:ipa_relation}. The extractor obtained this way runs in expected
  polynomial time once the rewinding oracle is folded into the query log.

  Now \(v_{1} \neq v_{2}\) forces \((a_{1}, r_{1}) \neq (a_{2}, r_{2})\), yet both
  reconstruct the \emph{same} \(P\). Subtracting the two Pedersen equations gives
  \[
    \inner{a_{1} - a_{2}}{g} + [r_{1} - r_{2}]H = O,
  \]
  a non-trivial discrete-log relation among the random generators \(g\), \(H\)
  (and \(U\), once the \([v]U\) terms are accounted for), because the coefficient
  vector \(a_{1} - a_{2}\) is nonzero on the monomial basis distinguishing
  \(v_{1}\) from \(v_{2}\). By the discrete log relation assumption (Halo
  Definition~6) such a relation occurs only with negligible probability, so both
  openings can succeed only with negligible probability. Hence the binding error
  \(\varepsilon_{\mathrm{b}}\) is the sum of the emulation gap of \cref{def:wee}
  and the discrete-log advantage, both negligible.

  The faithful construction of the rewinding emulator and the exact accounting of
  the extraction error are carried out in Halo Appendix~A and are not reproduced
  here; this chapter states the result and records the argument's shape.
\end{proof}
